%\documentclass[smaller,handout]{beamer}
% \documentclass[smaller,handout]{beamer}
\def\bmode{2} % Mode 0 for presentation, mode 1 for a handout with notes, mode 2 fo% r handout without notes
\if 0\bmode
\documentclass[smaller]{beamer}
\else \if 1\bmode
\immediate\write18{pdflatex -jobname=\jobname-Handout-Notes\space\jobname}
\documentclass[smaller,handout]{beamer}
\usepackage{handoutWithNotes}
\pgfpagesuselayout{2 on 1 with notes}[letterpaper, landscape, border shrink=4mm]
\else \if 2\bmode
\immediate\write18{pdflatex -jobname=\jobname-Handout\space\jobname}
\documentclass[smaller,handout]{beamer}
\fi
\fi
\fi

%%%%%%%%%%%%%%%%%%%%%%%%%%%%%%%%%%%%%%%%%%%%%%%%%%%%%%%%%%%%%%%%%%%%%%%%%%%%%%%%%%%%%%%%%%%%%
\newcommand{\coursetitle}{CEE 697M: Probabilistic Machine Learning}
\newcommand{\longlecturetitle}{M2 Linear Methods: Splines, GAMs and GLMs}
\newcommand{\shortlecturetitle}{L2e: Generalized Linear Models}
\newcommand{\instructor}{Jimi Oke}
\newcommand{\lecturedate}{October  7, 2025}
%%%%%%%%%%%%%%%%%%%%%%%%%%%%%%%%%%%%%%%%%%%%%%%%%%%%%%%%%%%%%%%%%%%%%%%%%%%%%%%%%%%%%%%%%%%%%

 
% \usepackage[T1]{fontenc} 
% \usepackage{lmodern} 
%\usepackage{etex}
 %\newcommand{\num}{6{} }

% \usetheme[
%   outer/progressbar=foot,
%   outer/numbering=fraction,
%   block=fill,
%   inner/subsectionpage=progressbar
% ]{metropolis}
\usetheme{Madrid}
\useoutertheme[subsection=false]{miniframes} % Alternatively: miniframes, infolines, split
\useinnertheme{circles}
% %\useoutertheme{Frankfurt}
% \usecolortheme{beaver}
% %\useoutertheme{crane}
% %\useoutertheme{metropolis}
\usepackage[backend=biber,style=authoryear,maxcitenames=2,maxbibnames=99,safeinputenc,url=false, eprint=false]{biblatex}
%\addbibresource{bib/references.bib}
% \AtEveryCitekey{\iffootnote{{\tiny}\tiny}{\tiny}}

% %\usepackage{pgfpages}
% %\setbeameroption{hide notes} % Only slides
% %\setbeameroption{show only notes} % Only notes
% %\setbeameroption{hide notes} % Only notes
% %\setbeameroption{show notes on second screen=right} % Both

% % \usepackage[sfdefault]{Fira Sans}

% % \setsansfont[BoldFont={Fira Sans}]{Fira Sans Light}
% % \setmonofont{Fira Mono}

% %\usepackage{fira}
% %\setsansfont{Fira}
% %\setmonofont{Fira Mono}
% % To give a presentation with the Skim reader (http://skim-app.sourceforge.net) on OSX so
% % that you see the notes on your laptop and the slides on the projector, do the following:
% % 
% % 1. Generate just the presentation (hide notes) and save to slides.pdf
% % 2. Generate onlt the notes (show only nodes) and save to notes.pdf
% % 3. With Skim open both slides.pdf and notes.pdf
% % 4. Click on slides.pdf to bring it to front.
% % 5. In Skim, under "View -> Presentation Option -> Synhcronized Noted Document"
% %    select notes.pdf.
% % 6. Now as you move around in slides.pdf the notes.pdf file will follow you.
% % 7. Arrange windows so that notes.pdf is in full screen mode on your laptop
% %    and slides.pdf is in presentation mode on the projector.

% % Give a slight yellow tint to the notes page
% \setbeamertemplate{note page}{\pagecolor{yellow!5}\insertnote}\usepackage{palatino}

% %\usetheme{metropolis}
% %\usecolortheme{beaver}
 \usepackage{tipa}
% \usepackage{enumerate}
\definecolor{darkcandyapplered}{HTML}{A40000}
\definecolor{lightcandyapplered}{HTML}{e74c3c}

% %\setbeamercolor{title}{fg=darkcandyapplered}

% \definecolor{UBCblue}{rgb}{0.04706, 0.13725, 0.26667} % UBC Blue (primary)
% \definecolor{UBCgrey}{rgb}{0.3686, 0.5255, 0.6235} % UBC Grey (secondary)

% % \setbeamercolor{palette primary}{bg=darkcandyapplered,fg=white}
% % \setbeamercolor{palette secondary}{bg=darkcandyapplered,fg=white}
% % \setbeamercolor{palette tertiary}{bg=darkcandyapplered,fg=white}
% % \setbeamercolor{palette quaternary}{bg=darkcandyapplered,fg=white}
% % \setbeamercolor{structure}{fg=darkcandyapplered} % itemize, enumerate, etc
% % \setbeamercolor{section in toc}{fg=darkcandyapplered} % TOC sections
% % \setbeamercolor{frametitle}{fg=darkcandyapplered,bg=white} % TOC sections
% % \setbeamercolor{title in head/foot}{bg=white,fg=white} % TOC sections
% % \setbeamercolor{button}{fg=darkcandyapplered} % TOC sections

% % % Override palette coloring with secondary
% % \setbeamercolor{subsection in head/foot}{bg=lightcandyapplered,fg=white}

%\usecolortheme{crane}
% \makeatletter
% \setbeamertemplate{headline}{%
%   \begin{beamercolorbox}[colsep=1.5pt]{upper separation line head}
%   \end{beamercolorbox}
%   \begin{beamercolorbox}{section in head/foot}
%     \vskip1pt\insertsectionnavigationhorizontal{\paperwidth}{}{}\vskip1pt
%   \end{beamercolorbox}%
%   \ifbeamer@theme@subsection%
%     \begin{beamercolorbox}[colsep=1.5pt]{middle separation line head}
%     \end{beamercolorbox}
%     \begin{beamercolorbox}[ht=2.5ex,dp=1.125ex,%
%       leftskip=.3cm,rightskip=.3cm plus1fil]{subsection in head/foot}
%       \usebeamerfont{subsection in head/foot}\insertsubsectionhead
%     \end{beamercolorbox}%
%   \fi%
%   \begin{beamercolorbox}[colsep=1.5pt]{lower separation line head}
%   \end{beamercolorbox}
% }
% \makeatother

% Reduce size of frame box
\setbeamertemplate{frametitle}{%
    \nointerlineskip%
    \begin{beamercolorbox}[wd=\paperwidth,ht=2.0ex,dp=0.6ex]{frametitle}
        \hspace*{1ex}\insertframetitle%
    \end{beamercolorbox}%
}


%\setbeamercolor{frametitle}{bg=darkcandyapplered!80!black!90!white}
%\setbeamertemplate{frametitle}{\bf\insertframetitle}

%\setbeamercolor{footnote mark}{fg=darkcandyapplered}
%\setbeamercolor{footnote}{fg=darkcandyapplered!70}
%\Raggedbottom
%\setbeamerfont{page number in head/foot}{size=\tiny}
%\usepackage[tracking]{microtype}


% %\usepackage[sc,osf]{mathpazo}   % With old-style figures and real smallcaps.
% %\linespread{1.025}              % Palatino leads a little more leading

% % Euler for math and numbers
% %\usepackage[euler-digits,small]{eulervm}
% %\AtBeginDocument{\renewcommand{\hbar}{\hslash}}
\usepackage{graphicx,multirow,booktabs}
\usepackage{animate}
\usepackage{media9}


% %\mode<presentation> { \setbeamercovered{transparent} }

\setbeamertemplate{navigation symbols}{}
\makeatletter
\def\beamerorig@set@color{%
  \pdfliteral{\current@color}%
  \aftergroup\reset@color
}
\def\beamerorig@reset@color{\pdfliteral{\current@color}}
\makeatother


% %=== GRAPHICS PATH ===========
\graphicspath{{./m2-images/}}
% % Marginpar width
% %Marginpar width
% %\setlength{\marginparsep}{.02in}


% %% Captions
% % \usepackage{caption}
% % \captionsetup{
% %   labelsep=quad,
% %   justification=raggedright,
% %   labelfont=sc
% % }

% \setbeamerfont{caption}{size=\footnotesize}
% \setbeamercolor{caption name}{fg=darkcandyapplered}

% %AMS-TeX packages

\usepackage{amssymb,amsmath,amsthm,mathtools} 
\usepackage{bm}
\DeclareMathOperator*{\argmax}{arg\,max}
\DeclareMathOperator*{\argmin}{arg\,min}
% \usepackage{color}

% %https://tex.stackexchange.com/a/31370/2269
\usepackage{mathtools,cancel}

\renewcommand{\CancelColor}{\color{red}} %change cancel color to red

\makeatletter
\let\my@cancelto\cancelto %copy over the original cancelto command
\newcommand<>{\cancelto}[2]{\alt#3{\my@cancelto{#1}{#2}}{\mathrlap{#2}\phantom{\my@cancelto{#1}{#2}}}}
% redefine the cancelto command, using \phantom to assure that the
% result doesn't wiggle up and down with and without the arrow
\makeatother


% %\usepackage{comment}
% %\usepackage{hyperref,enumerate}
% \usepackage{minitoc,array}

% \definecolor{slblue}{rgb}{0,.3,.62}
% % \hypersetup{
% %     colorlinks,%
% %     citecolor=blue,%
% %     filecolor=blue,%
% %     linkcolor=blue,
% %     urlcolor=slblue
% % }

% \usepackage{epstopdf}
% \epstopdfDeclareGraphicsRule{.gif}{png}{.png}{convert gif:#1 png:\OutputFile}
% \AppendGraphicsExtensions{.gif}

% %\usepackage{listings}

% %%% TIKZ
% \usepackage{forest}
\usepackage{tikz}
\usepackage{pgfplots}
\usepackage{pgfplotstable}
%\usepackage{pgfgantt}
\pgfplotsset{compat=newest}

\usetikzlibrary{fit,arrows,shapes,positioning,shapes.geometric}
\usetikzlibrary{decorations.markings}
\usetikzlibrary{shadows,automata}
\usetikzlibrary{patterns}
\usetikzlibrary{trees,mindmap,backgrounds}
%\usetikzlibrary{circuits.ee.IEC}
\usetikzlibrary{decorations.text}
% % For Sagnac Picture
% \usetikzlibrary{%
%     decorations.pathreplacing,%
%     decorations.pathmorphing%
% }
% \tikzset{no shadows/.style={general shadow/.style=}}
% %
% %\usepackage{paralist}

% \tikzset{
%   font=\Large\sffamily\bfseries,
%   red arrow/.style={
%     midway,red,sloped,fill, minimum height=3cm, single arrow, single arrow head extend=.5cm, single arrow head indent=.25cm,xscale=0.3,yscale=0.15,
%     allow upside down
%   },
%   black arrow/.style 2 args={-stealth, shorten >=#1, shorten <=#2},
%   black arrow/.default={1mm}{1mm},
%   tree box/.style={draw, rounded corners, inner sep=1em},
%   node box/.style={white, draw=black, text=black, rectangle, rounded corners},
% }

% %%% FORMAT PYTHON CODE
% %\usepackage{listings}
% % Default fixed font does not support bold face
% \DeclareFixedFont{\ttb}{T1}{txtt}{bx}{n}{8} % for bold
% \DeclareFixedFont{\ttm}{T1}{txtt}{m}{n}{8}  % for normal

% % Custom colors
% \definecolor{deepblue}{rgb}{0,0,0.5}
% \definecolor{deepred}{rgb}{0.6,0,0}
% \definecolor{deepgreen}{rgb}{0,0.5,0}

% %\usepackage{animate}

% % Python style for highlighting
% % \newcommand\pythonstyle{\lstset{
% % language=Python,
% % basicstyle=\footnotesize\ttm,
% % otherkeywords={self},             % Add keywords here
% % keywordstyle=\footnotesize\ttb\color{deepblue},
% % emph={MyClass,__init__},          % Custom highlighting
% % emphstyle=\footnotesize\ttb\color{deepred},    % Custom highlighting style
% % stringstyle=\color{deepgreen},
% % frame=tb,                         % Any extra options here
%     % showstringspaces=false            % 
% % }}

% % % Python environment
% % \lstnewenvironment{python}[1][]
% % {
% % \pythonstyle
% % \lstset{#1}
% % }
% % {}

% % % Python for external files
% % \newcommand\pythonexternal[2][]{{
% % \pythonstyle
% % \lstinputlisting[#1]{#2}}}

% % Python for inline
% % 
% % \newcommand\pythoninline[1]{{\pythonstyle\lstinline!#1!}}

% %\usepackage{algorithm2e}

\newcommand{\eps}{\epsilon}
\newcommand{\bX}{\mb X}
\newcommand{\by}{\mb y}
\newcommand{\bbe}{\bm\beta}
\newcommand{\beps}{\bm\epsilon}
\newcommand{\bY}{\mb Y}

\newcommand{\osn}{\oldstylenums}
\newcommand{\dg}{^{\circ}}
\newcommand{\lt}{\left}
\newcommand{\rt}{\right}
\newcommand{\pt}{\phantom}
\newcommand{\tf}{\therefore}
\newcommand{\?}{\stackrel{?}{=}}
\newcommand{\fr}{\frac}
\newcommand{\dfr}{\dfrac}
\newcommand{\ul}{\underline}
\newcommand{\tn}{\tabularnewline}
\newcommand{\nl}{\newline}
\newcommand\relph[1]{\mathrel{\phantom{#1}}}
\newcommand{\cm}{\checkmark}
\newcommand{\ol}{\overline}
\newcommand{\rd}{\color{red}}
\newcommand{\bl}{\color{blue}}
\newcommand{\pl}{\color{purple}}
\newcommand{\og}{\color{orange!90!black}}
\newcommand{\gr}{\color{green!40!black}}
\newcommand{\dca}{\color{darkcandyapplered}}
\newcommand{\nin}{\noindent}
\newcommand*\circled[1]{\tikz[baseline=(char.base)]{
            \node[shape=circle,draw,thick,inner sep=1pt] (char) {\small #1};}}

\newcommand{\bc}{\begin{compactenum}[\quad--]}
\newcommand{\ec}{\end{compactenum}}

\newcommand{\p}{\partial}
\newcommand{\pd}[2]{\frac{\partial{#1}}{\partial{#2}}}
\newcommand{\dpd}[2]{\dfrac{\partial{#1}}{\partial{#2}}}
\newcommand{\pdd}[2]{\frac{\partial^2{#1}}{\partial{#2}^2}}
\newcommand{\pde}[3]{\frac{\partial^2{#1}}{\partial{#2}\partial{#3}}}
\newcommand{\nmfr}[3]{\Phi\left(\frac{{#1} - {#2}}{#3}\right)}
\newcommand{\Err}{\text{Err}}
\newcommand{\err}{\text{err}}

%\DeclarePairedDelimiter\ceil{\lceil}{\rceil}
%\DeclarePairedDelimiter\floor{\lfloor}{\rfloor}

%%%% GREEK LETTER SHORTCUTS %%%%%
\newcommand{\la}{\lambda}
\renewcommand{\th}{\theta}
\newcommand{\al}{\alpha}
\newcommand{\G}{\Gamma}
\newcommand{\si}{\sigma}
\newcommand{\Si}{\Sigma}


\pgfmathdeclarefunction{poiss}{1}{%
  \pgfmathparse{(#1^x)*exp(-#1)/(x!)}%
  }

\pgfmathdeclarefunction{gauss}{2}{%
  \pgfmathparse{1/(#2*sqrt(2*pi))*exp(-((x-#1)^2)/(2*#2^2))}%
}

\pgfmathdeclarefunction{expo}{2}{%
  \pgfmathparse{#1*exp(-#1*#2)}%
}

\pgfmathdeclarefunction{expocdf}{2}{%
  \pgfmathparse{1 -exp(-#1*#2)}%
}

\newcommand{\mb}{\mathbb}
\newcommand{\mc}{\mathcal}
\newcommand{\tr}{^{\top}}
\newcommand{\pe}{\pause}
% \usepackage{pst-plot}

% \usepackage{pstricks-add}
% \usepackage{auto-pst-pdf}   

% \psset{unit = 3}

% \def\target(#1,#2){%
%  {\psset{fillstyle = solid}
%   \rput(#1,#2){%
%     \pscircle[fillcolor = white](0.7,0.7){0.7}
%     \pscircle[fillcolor = blue!60](0.7,0.7){0.5}
%     \pscircle[fillcolor = white](0.7,0.7){0.3}
%     \pscircle[fillcolor = red!80](0.7,0.7){0.1}}}}
% \def\dots[#1](#2,#3){%
%     \psRandom[
%       dotsize = 2pt,
%       randomPoints = 25
%     ](!#2 #1 0.04 sub sub #3 #1 0.04 sub sub)%
%      (!#2 #1 0.04 sub add #3 #1 0.04 sub add)%
%      {\pscircle[linestyle = none](#2,#3){#1}}}


%%%%%%%%%%%%%%%%%%%%%%%%%%%%%%%%%%%%%%%%%%%%%%%%%%%
%%%%%%%%%%%%%%%%%%%%%%%%%%%%%%%%%%%%%%%%%%%%%%%%%%%
\title[\shortlecturetitle]{ {\normalsize \coursetitle}
  \\ \longlecturetitle}
\date[\lecturedate]{\footnotesize \lecturedate}
\author{{\bf \instructor}}
\institute[UMass Amherst]{
%\titlegraphic{\hfill
  \begin{tikzpicture}[baseline=(current bounding box.center)]
    \node[anchor=base] at (-7,0) (its) {\includegraphics[scale=.3]{UMassEngineering_vert}} ;
  \end{tikzpicture}
  % \hfill\includegraphics[height=1.5cm]{logo}
}

%https://tex.stackexchange.com/questions/55806/mindmap-tikzpicture-in-beamer-reveal-step-by-step
  \tikzset{
    invisible/.style={opacity=0},
    visible on/.style={alt={#1{}{invisible}}},
    alt/.code args={<#1>#2#3}{%
      \alt<#1>{\pgfkeysalso{#2}}{\pgfkeysalso{#3}} % \pgfkeysalso doesn't change the path
    },
  }


% https://tex.stackexchange.com/questions/446468/labels-with-arrows-for-an-equation
% https://tex.stackexchange.com/a/402466/121799
\newcommand{\tikzmark}[3][]{
\ifmmode
\tikz[remember picture,baseline=(#2.base)] \node [inner sep=0pt,#1](#2) {$#3$};
\else
\tikz[remember picture,baseline=(#2.base)] \node [inner sep=0pt,#1](#2) {#3};
\fi
}

% \lstset{language=matlab,
%                 basicstyle=\scriptsize\ttfamily,
%                 keywordstyle=\color{blue}\ttfamily,
%                 stringstyle=\color{blue}\ttfamily,
%                 commentstyle=\color{gray}\ttfamily,
%                 morecomment=[l][\color{gray}]{\#}
%               }


              
\begin{document}

\maketitle

\begin{frame}
  \frametitle{Outline}
  \tableofcontents
\end{frame}


 
 


\section{Backgrounnd: Exponential family}
  \begin{frame}
    \frametitle{The exponential family}
    \pause

    A probability distribution belongs to the exponential family if its density can be modeled as:

    \begin{equation}
      p(\bm y |\bm \eta) =   \fr{1}{Z(\bm \eta)} h(\bm y) \exp\lt[ \bm \eta\tr\mathcal{T} (\bm y)\rt]
      \pause = 
       h(\bm y) \exp\lt[\bm \eta\tr\mathcal{T} (\bm y) - A(\bm\eta)\rt]
    \end{equation}
    \pause

    where:
    \begin{itemize}
      \item $Z(\bm \eta)$ is the partition function (normalization constant) \pause
      \item $h(\bm y)$ is the base measure (scaling constant; typically 1)\pause
      \item $\bm \eta$ are the natural/canonical parameters \pause
      \item $\mathcal{T}(\bm y)$ are the sufficient statistics \pause
      \item $A(\bm \eta) = \ln Z(\bm \eta)$ is the log-partition function \pause
    \end{itemize}
  \pause

  The log-likelihood is then given by:
  \pause
  \begin{equation}
    \log p(\bm y |\bm \eta) = \log h(\bm y) + \bm \eta\tr\mathcal{T} (\bm y) - A(\bm\eta) + \text{const}
  \end{equation}
  \end{frame}

  \begin{frame}
    \frametitle{Properties of exponential family}\pause

\begin{itemize}
  \item Generalization: 
  we define $\bm\eta = f(\phi)$, thus:
  \begin{equation}
    p (\bm y |\bm\phi) = h(\bm y) \exp\lt[f(\bm \phi)\tr\mathcal{T} (\bm y) - A(f(\bm\phi))\rt]
  \end{equation}
  \pause
  \item If $f(\bm\phi)$ is nonlinear, then the model is in the curved exponential family\pause
  \item If $\bm\eta = f(\bm\phi) = \bm\phi$, the model is in \textbf{canonical form} \pause 
  \item If $\mathcal{T}(\bm y) = \bm y$, the model is in the natural exponential family \pause
  \begin{equation}
    p(\bm y|\bm \eta) = h (\bm y) \exp\lt[\bm \eta\tr \bm y - A(\bm \eta)\rt]
  \end{equation}
\end{itemize}  
  
  \end{frame}

\begin{frame}
  \frametitle{Bernoulli distribution in exponential family form (1/2)}
\pause

  The Bernoulli distribution is given by:
  \begin{equation}
    p(y|\mu) = \mu^{y}(1 - \mu)^{1 - y}, \quad y \in \{0,1\}, \quad 0 < \mu < 1
  \end{equation}
  \pause
  where $\mu = \mathbb{E}(y)$ is the probability of success. \pause Rewriting: \pause
  \begin{eqnarray*}
    p(y|\mu) &=& (1 - \mu) \lt(\fr{\mu}{1 - \mu}\rt)^{y} \pause
             =\pause (1 - \mu) \exp\lt[y\log\lt(\fr{\mu}{1 - \mu}\rt)\rt] \\ \pause
             &=& {\gr (1 - \mu)} \exp\lt[{\pl y} {\rd \log\lt(\fr{\mu}{1 - \mu}\rt)} - {\bl 0}\rt]
  \end{eqnarray*}
  \pause

  Comparing to the exponential family form:
  \begin{eqnarray*}
    \gr h(y) &=& \gr 1 - \mu  \quad \pause \text{(base measure)}\\ \pause
    \pl \mathcal{T}(y) &=& \pl y \quad \pause \text{(sufficient statistic)}\\ \pause
    \rd \eta &=& \rd \log\lt(\fr{\mu}{1 - \mu}\rt) \quad \pause \text{(natural parameter)} \\ \pause
    \bl A(\eta) &=& \bl 0 \quad \pause \text{(log-partition function)}
  \end{eqnarray*}
  

\end{frame}

\begin{frame}
  \frametitle{Cumulant generating function}
\pause
\begin{itemize}
  \item Cumulants $\kappa_n(\bm y)$ are functions of the central moments of a distribution \pause
  \item For example, $\kappa_1(\bm y) = \mathbb{E}(\bm y)$ and $\kappa_2(\bm y) = \mathbb{V}(\bm y)$ \pause
  \item Higher order cumulants are polynomial functions of the central moments \pause
  \item The cumulants of a distribution are defined by the cumulant generating function (CGF):
    \begin{equation}
      K_{\bm y}(t) = \log \mathbb{E}(\exp(t\bm y))
    \end{equation}
    \pause
    where $\mathbb{E}(\exp(t\bm y))$ is the moment generating function (MGF) of $\bm x$ \pause

\item In the exponential family, the log-partition function $A(\bm\eta)$ is the CGF of the sufficient statistics $\mathcal{T}(\bm y)$ \pause
  \item Thus, the cumulants can be obtained by differentiating $A(\bm\eta)$:
    \begin{eqnarray*}
      \kappa_1(\mathcal{T}(\bm y)) &=& \mathbb{E}(\mathcal{T}(\bm y)) = \nabla_{\bm\eta} A(\bm\eta) \\ \pause
      \kappa_2(\mathcal{T}(\bm y)) &=& \text{Cov}(\mathcal{T}(\bm y)) = \nabla^2_{\bm\eta} A(\bm\eta)
    \end{eqnarray*}
\end{itemize}
\end{frame}

\begin{frame}
  \frametitle{Unique global maximum of the likelihood}\pause

  From the CGF properties, we have: 
  \begin{equation}
    \nabla^2_{\bm\eta} A(\bm\eta) = \text{Cov}(\mathcal{T}(\bm y)) > 0
  \end{equation}
  This implies that the log-partition function $A(\bm\eta)$ is strictly convex. \pause Thus, the log-likehood
  \pause

  \begin{equation}
    \log p(\bm y |\bm \eta) = \log h(\bm y) + \bm \eta\tr\mathcal{T} (\bm y) - A(\bm\eta) + \text{const}
  \end{equation}
  
  is guaranteed to have a unique global maximum.

\end{frame}
\section{GLMs}

\begin{frame}
  \frametitle{The generalized linear model (GLM)}
  \label{glm}
  \pause

  \begin{itemize}
  \item Conventional linear regression models have the form:
    \begin{equation}
      p(y|\bm x, \bm w) \sim \mathcal{N}(y| \bm x\tr \bm w, \sigma^{2})
    \end{equation}
    \pause where \pause
    \begin{itemize}
     \item $y_{n}$ is a continuous response
     \item $\bm x_{n}$ is a vector of quantitative and/or qualitative explanatory variables
     \end{itemize}
     \pause

    \item Generalized linear models (GLMs) were introduced to extend this framework to allow $y_{n}$ to be modeled by other \href{https://people.eecs.berkeley.edu/~jordan/courses/260-spring10/other-readings/chapter8.pdf}{exponential family distributions} besides the normal/Gaussian, e.g.
     \begin{itemize}
     \item exponential
     \item binomial/multinomial (with fixed number of trials)
     \item Poisson
     \end{itemize}
    \pause

  \item In the GLM framework: \pause
    \begin{itemize}
    \item The mean of $y_{n}$ is given by $\mu_{n}$
    \item $\mu_{n}$ can be specified by a nonlinear function of $\bm x_{n}\tr \bm w$
    \item Note that the simple linear regression is a special case of GLM in which 
    $\mu_{n} = \bm x_{n}\tr \bm w$ and $y_{n}$ follows a Gaussian distribution
    \end{itemize}
  \end{itemize}
  %  \pause

   % \hyperlink{recap}{\beamerbutton{Back to Recap}}
  \end{frame}

  \begin{frame}
    \frametitle{GLM formulation}\pause

    The GLM is a version of the exponential family distribution in which the natural parameters $\eta_n$ are a {\bf linear function} of the output.\footnote{The technical name of the GLM form of the distribution is the \textbf{exponential dispersion model/family, often abbreviated as ``EDM.''}} \pause It is given by:\pause

    \begin{equation}
      p(y_n |\bm x_n, \bm w, \sigma^2) = \exp\lt[ \fr{y_n\eta_n - A(\eta_n)}{\sigma^2} + \log h(y_n,\sigma^2) \rt]
    \end{equation}
    \pause
    where:
    \pause
    \begin{itemize}
    \item $\eta_n = \bm w \tr \bm x_n$ is the natural parameter (input) \pause
    %\item $\mu_n = E(y_n) = g^{-1}(\eta_n)$ is the mean of $y_n$ \pause
    %\item $g(\cdot)$ is the link function \pause
    \item $y_n  = \mathcal{T}(y_n)$ is the sufficient statistic \pause
    \item $A(\eta_n)$ is the log-partition function (or log normalizer)\pause
    \item $h(y_n,\sigma^  2)$ is the base measure \pause
    \item $\sigma^2$ is the {\bf dispersion parameter}  (typically known or set to 1)
    \end{itemize}
  \end{frame}

  \begin{frame}
    \frametitle{Link and mean functions}\pause



    Recalling that the mean and variance of the sufficient statistics $\mathcal{T}(y_n) = y_n$ are given by the first and second derivatives of the log-partition function $A(\eta_n)$, we have:
    \pause
  \begin{eqnarray}
    \mathbb{E}(y_n |\bm x_n ,\bm w, \sigma^2) &=& A'(\eta_n) = \pause \ell^{-1}(\eta_n)  \\\pause
    \mathbb{V}(y_n|\bm x_n ,\bm w, \sigma^2) &=& A''(\eta_n)\sigma^2 \pause
  \end{eqnarray}
    
  \pause
  We define the \textbf{mean function} as 
  \begin{equation}\mu_n = \ell^{-1}(\eta_n) \end{equation} \pause 
  and the \textbf{link function} as its inverse: 
  \begin{equation}
    g(\mu_n) = \ell(\mu_n)  
  \end{equation}
  \pause The link function\footnote{In most textbooks, the link function is denoted as $g(\mu)$, while Murphy uses $\ell(\mu)$.} is thus the inverse of the mean function, and its role is to map the mean output/response $\mu_n$ to the linear predictor $\eta_n = \bm w \tr \bm x_n$.
  \end{frame}

  \begin{frame}
    \frametitle{Linear regression (1/2)}

    \pause

    Linear regression has the form:

    \begin{equation}
      p(y_n |\bm x_n, \bm w, \sigma^2) = \fr{1}{\sqrt{2\pi\sigma^2}} \exp\lt(-\fr{1}{2\sigma^2}(y_n - \bm w \tr \bm x_n)^2\rt)
    \end{equation}
    \pause
    Taking logs: \pause
    \begin{equation}
      \log p(y_n |\bm x_n, \bm w, \sigma^2) = -\fr{1}{2}\log(2\pi\sigma^2) - \fr{1}{2\sigma^2}(y_n - \bm w \tr \bm x_n)^2
    \end{equation}\pause

    Setting $\eta_n = \bm w \tr \bm x_n$, we can write in GLM form as:
    \begin{equation}
        \log p(y_n|\bm x_n, \bm w, \sigma^2) = \fr{y_n \eta_n - \eta_n^2/2}{\sigma^2} - \fr12 \lt( \fr{y_n^2}{\sigma^2} + \log(2\pi\sigma^2) \rt)
    \end{equation}
    \pause
  \end{frame}

  \begin{frame}
    \frametitle{Linear regression (2/2)}
      If we set: \pause 
    \begin{eqnarray}
       A(\eta_n) &=& \eta_n^2/2 \\\pause
       h(y_n,\sigma^2) &=& \fr{1}{\sqrt{2\pi\sigma^2}} \exp\lt(-\fr{1}{2\sigma^2}y_n^2\rt)
    \end{eqnarray}\pause
   then we can write:\pause
    \begin{equation}
     \log p(y_n |\bm x_n, \bm w, \sigma^2) = \fr{y_n \eta_n - A(\eta_n)}{\sigma^2} + \log h(y_n,\sigma^2)
    \end{equation}
    \pause
    And thus, the cumulants are given by:
    \begin{eqnarray}
      \mathbb{E}(y_n |\bm x_n ,\bm w, \sigma^2) &=& A'(\eta_n) \pause = \eta_n = \bm w \tr \bm x_n \\
      \text{Var}(y_n|\bm x_n ,\bm w, \sigma^2) &=& A''(\eta_n)\sigma^2 = \sigma^2
    \end{eqnarray}
    \pause
    Thus, the mean function is $\mu_n = \ell^{-1}(\eta_n) = \eta_n$ and the link function is $g(\mu_n) = \ell(\mu_n) = \mu_n$ (identity link function).
  \end{frame}

  \section{Link functions}
  \begin{frame}
    \frametitle{GLM components}

    \pause

    A GLM can be considered as consisting of three parts: \pause

    \begin{itemize}
    \item \textbf{Random component:} this is the probability distribution of the response variable $p(\bm y|\bm X, \bm w)$\pause
    \item \textbf{Systematic component:} specifies the explanatory variables within the linear combination of their coefficients ($\bm X\bm w $)\pause
    \item \textbf{Link function $\bm{g(\mu)}$:} defines the relationship between the random and systematic components: \pause
      \begin{itemize}
      \item Simple linear regression ({\bf identity} link function): \pause
        \begin{equation}
          g(\mu_{n}) = g(\mathbb{E}(y_{n })) = \bm w\tr \bm x_n
        \end{equation}
        \pause
      \item Binary logistic regression ({\bf logit} link function): \pause
        \begin{equation}
          g(\mu_{n}) = g(p(\bm x_{n})) = \pause \text{logit} (p(\bm x_{n})) = \pause
          \ln\lt(\fr{p(\bm x_{n})}{1 - p(\bm x_{n})} \rt) = \pause \bm w\tr \bm x_n
        \end{equation}
      \end{itemize}
    \end{itemize}
  \end{frame}



  \begin{frame}
    \frametitle{Assumptions of GLM}
    \pause
    \begin{itemize}
    \item The observations of the response variable $\bm y$  are i.i.d.
      \pause
    \item Response variable $y_{n}$ is typically exponentially distributed (not restricted to being normally distributed) \pause
      \begin{itemize}
      \item Implies that errors need not be normally distributed (but should be independent)
      \end{itemize}
    \item Link function ($g(\mu_n)$) is linear with respect to the coefficients ($ w _{d}$)\pause
      \begin{itemize}
      \item Relationship between response and explanatory variables does not have to be linear \pause
      \item Explanatory variables can be nonlinear transformations of original values (as in simple linear regression)
      \end{itemize}
      \pause
    \item Variance may not homogeneous (i.e.\ homoscedasticity is not a requirement)\pause
    \item Parameters are estimated via MLE 
    \end{itemize}
  \end{frame}

  \begin{frame}
    \frametitle{Commonly used GLM models and their components}

    \pause
    \begin{center}\scriptsize
      \begin{tabular}{l p{.6in} p{.5in} l }\toprule
      \bf Model & \bf Random component& \bf Mean/output ($\mu_n$) & \bf Link function   \\\midrule
      Linear regression & Gaussian & $\bm w \tr \bm x_{n}$ &  Identity: $g(\mu_n) = \mu_{n} =  \bm w \tr \bm x_{n}$ \\\midrule
      Binary logistic regression & Bernoulli  & $\bm\sigma(\bm w \tr \bm x_{n})$ & Logit:  $g(\mu_{n}) = \log\lt(\fr{\mu_{n}}{1 - \mu_{n}}\rt)$ \\\midrule
      Probit regression & Bernoulli & $\bm\sigma(\bm w \tr \bm x_{n})$ & Probit: $g(\mu_{n}) = \Phi^{-1}(\mu_{n})$ \\\midrule
      Multinomial logit/logistic  & Categorical & $\mathcal{S}(\bm W \bm x_{n})$ & Multinomial logit: $g(\mu_{nc}) = \log\lt(\fr{\mu_{nc}}{\mu_{nC}}\rt)$ \\\midrule
      Poisson regression & Poisson & $\exp(\bm w \tr \bm x_{n})$ & Log: $g(\mu_{n}) = \log(\mu_{n})$ \\\bottomrule
      \end{tabular}
    \end{center}

    \pause

    Note that in all cases, the link function always results in:
    \begin{equation}
      g(\mu_n) = \bm w \tr \bm x_{n} = \eta_n
    \end{equation}
    \pause Its job is to ``link'' the response to the systematic component via a suitable transformation that results in a linear function of the $ w $'s.
  \end{frame}


  \begin{frame}
    \frametitle{Canonical and non-canonical link functions}

    \pause

    Link functions can be classified as either canonical or non-canonical: \pause
    \begin{itemize}
    \item Canonical link function: results in the  canonical/natural parameters $\theta$ of the random component (i.e. $\eta = g(\mu) = \theta$) \pause
      \begin{itemize}
      \item \textbf{identity} (linear regression): $g(\mu) = \mu$; $\mu$ is the canonical parameter of the Gaussian distribution
      \item  \textbf{logit} (binary logistic regression): $g(\mu) = \log\lt(\fr{\mu}{1 - \mu}\rt)$; the logit is the canonical parameter of the Bernoulli distribution
      \item \textbf{log} (Poisson regression): $g(\mu) = \log(\mu)$; $\log(\mu)$ is the canonical parameter of the Poisson distribution
      \end{itemize}
      \pause
    \item Non-canonical link function: does not result in natural parameter of the random component/underlying distribution (i.e. $g(\mu_) \neq \mu_n$). \pause Examples:
    \begin{itemize}
      \item \textbf{probit} (binary probit regression): $g(\mu) = \Phi^{-1}(\mu)$; the probit is not the canonical parameter of the Bernoulli distribution
      \item \textbf{complementary log-log} (binary regression): $g(\mu) = \log(-\log(1 - \mu))$; the complementary log-log is not the canonical parameter of the Bernoulli distribution
    \end{itemize}
    \end{itemize}
  \end{frame}

  \begin{frame}
    \frametitle{Canonical link function (another view)}\pause

    Let $\theta$ be the natural parameter, $\eta$ the linear predictor, and $\mu$ the mean of the response variable. \pause
    Recall:

    \begin{eqnarray}
      \mu &=& \mathbb{E}(y|\theta) = A'(\theta) \pause \\
      \eta &=& \bm w \tr \bm x = g(\mu) = g(A'(\theta)) 
    \end{eqnarray}
    \pause

    If the link function is canonical, then:
    \begin{equation}
      g(\mu) =  (A')^{-1}(\mu) = \theta
    \end{equation}

    If not:
    \begin{equation}
      g(\mu) = (A')^{-1}(\mu) \neq \theta
    \end{equation}

    
  
  \end{frame}

  
  \begin{frame}
    \frametitle{Usage of canonical vs non-canonical link functions}

    \pause

    \begin{itemize}
      \item Canonical link functions are simpler to estimate and have desirable statistical properties. \pause
    \item Non-canonical link functions may be used when the canonical link does not provide a good fit to the data or when interpretability of the model is a priority. \pause
    \item In certain cases (e.g. probit regression), the non-canonical link (inverse CDF) allows for more efficient computation of the likelihood (e.g. Gibbs sampling)
    \end{itemize}
    

  \end{frame}

  \section{Fitting a GLM}

  \begin{frame}
    \frametitle{Fitting a GLM}

    \pause

    The process of fitting a GLM involves the following steps:

    \begin{itemize}
    \item Specify the distribution of the response variable (e.g. Gaussian, Bernoulli, Poisson)
    \item Choose a link function (canonical or non-canonical)
    \item Estimate the model parameters (e.g. using MLE)
    \item Assess the model fit (e.g. using residuals, AIC/BIC)
    \end{itemize}

  \end{frame}

\begin{frame}
  \frametitle{MLE of GLM parameters}
\pause
The negative log-likelihood (ignoring constant terms) is given by\pause


\begin{equation}
  \text{NLL}(\bm w) = -\log p(\mathcal{D}|\bm w) = -\sum_{n=1}^{N} \log p(y_n |\bm x_n, \bm w) = \sum_{n=1}^{N} \fr{A(\eta_n)}{\sigma^2} - \fr{y_n\eta_n}{\sigma^2}
\end{equation}
\pause

If we set $\mathcal{L}_n = \eta_n y_n - A(\eta_n)$, then the NLL can be written as:
\pause
\begin{equation}
  \text{NLL}(\bm w) = -\sum_{n=1}^{N} \fr{\mathcal{L}_n}{\sigma^2}
\end{equation}
\pause
where $\eta_n = \bm w \tr \bm x_n$. \pause

The gradient of the NLL (for a single term) is then given by:
\pause

\begin{equation}
  \bm g_n  =   \fr{y_n -\mu_n }{\sigma^2} \bm x_n
\end{equation}
\pause
where $\mu_n = A'(\eta_n) = \ell^{-1}(\eta_n)$ is the mean function.

\end{frame}

\section{Outlook}
\begin{frame}
  \frametitle{Summary}

  \begin{itemize}
    \item In the exponential family, the log-partition function $A(\theta)$ is the cumulant generating function of the sufficient statistics $\mathcal{T}(\bm y)$ \pause
    \item The log-partition function is strictly convex, thus the likelihood has a unique global maximum \pause
    \item GLMs are a flexible class of models that extend linear regression to handle non-normal response variables \pause
    \item Link functions $g(\mu)$ map the mean of the response variable to the linear predictor $\eta=\bm w\tr \bm x$ \pause
    \item Canonical link functions result in the natural parameters $\theta$ of the underlying distribution, while non-canonical link functions do not \pause
    \item GLM parameters $\eta$ can be estimated via MLE using gradient-based optimization methods
  \end{itemize}
\end{frame}
\end{document}

%%% Local Variables:
%%% mode: latex
%%% TeX-master: t
%%% End:
